\scanpage{46}\mspage{39}
% archive sheet number: 45
\[
 \alpha_{s}^{*}(\unsure{\Phi|\mathscr P}(I)_{*})\subset\Phi_{*}.
\]
b) Supposons que les $\operatorname{Hom}_{I}(t,s)$ soient finis (p. ex. $I$ finie). Alors \deleted{\unsure{les}} l'hyp. de a) s'applique pour $\Phi\unsure{=|\in}C$ et pour $\Phi=C\cap W$, \deleted{donc} \deleted{\unsure{et la structure}} la \deleted{structure} condition (iii) p \unsure{35|36} est satisfaite pour la structure $(\inserted{W(I),}C(I),TC(I),F(I),TF(I))$ dans $M(I)$. Même conclusion sans hyp. de finitude sur les \deleted{$M(I)$,} \inserted{$\operatorname{Hom}_{I}(t,s)$,} quand on suppose $M$ stable par ``petites'' sommes directes, et $C$ et $\deleted{\unsure{TC}}W$ stables également par de telles sommes directes.

\uline{NB} Vu le \unsure{Scholie|scholie} plus haut, il est parfaitement \unsure{raisonnable|vraisemblable} de se borner \unsure{au besoin|en théorie} aux catégories $I$ \uline{finies} -- vu que \deleted{\unsure{les}} \unsure{catégorie} de ces $I$ est stable par les produits $I\unsure{\mapsto|\vdash}I\times\unsure{\Psi|\Upsilon}$, $I\unsure{\mapsto|\vdash}I\times\Phi$.

Il reste, pour être heureux, à vérifier la condition de factorisation \textcircled{\scriptsize i} pour $(C,TF),(TC,F)$ certes, mais aussi et, surtout pour les $(C(I),TF(I)),(TC(I),F(I))$. C'est dire que de partir \deleted{\unsure{d'une}} d'une structure de Quillen close (pays de cocagne)
